\section{Representasi Data Finansial}

\subsection{Analisis Kualitatif dalam Pasar Finansial}
Analisis kualitatif atau yang lebih dikenal sebagai analisis fundamental dalam pasar finansial merupakan proses ekstraksi kondisi makro dan mikro secara mendalam yang menjadi fondasi utama di balik dinamika pergerakan harga aset dalam jangka panjang. Pendekatan fundamental ini secara spesifik berfokus pada berbagai variabel krusial seperti tingkat kesehatan industri, posisi strategis pasar, serta estimasi nilai intrinsik yang dimiliki oleh suatu entitas bisnis tertentu. Tingkat efisiensi pasar sangat bergantung pada kecepatan sistem dalam menyerap berita fundamental ini, sehingga informasi kualitatif menjadi penentu utama dalam pembentukan harga yang adil bagi seluruh pelaku pasar \cite{fama_1970}. Interaksi yang kompleks antara berbagai data fundamental tersebut pada akhirnya akan menentukan konsensus pasar terhadap ekspektasi risiko dan imbal hasil yang diharapkan oleh para agen secara kolektif. Pemahaman yang komprehensif terhadap aspek kualitatif ini memungkinkan investor untuk mengidentifikasi peluang investasi yang memiliki basis nilai riil yang kuat di tengah volatilitas pasar.

Dalam praktiknya, kapitalisasi pasar (\textit{market capitalization}) berfungsi sebagai indikator kualitatif yang esensial untuk mengukur skala likuiditas serta tingkat kepercayaan publik terhadap profil risiko suatu entitas finansial di pasar global. Laporan keuangan (\textit{financial statements}) kemudian memberikan rincian performa laba dan beban yang sangat mendetail, di mana data tersebut menjadi parameter utama dalam menentukan derajat kesehatan fiskal suatu perusahaan secara periodik. Variabel modern lainnya seperti \textit{fully diluted valuation} (FDV) dan analisis sentimen berita turut memberikan gambaran mengenai potensi saturasi nilai aset di masa depan yang sering kali sulit diprediksi secara linear. Pemahaman mendalam terhadap rincian data kualitatif ini menjadi syarat mutlak untuk membangun model ekonofisika yang mampu menangkap kompleksitas sistem keuangan secara utuh dan representatif. Integrasi data kualitatif ke dalam model matematis membantu mereduksi bias subjektif dalam pengambilan keputusan investasi yang bersifat strategis dan jangka panjang.

\subsection{Representasi Data pada Pergerakan Harga}
Analisis teknikal secara fundamental berfokus pada pemetaan stokastik aktivitas transaksi ke dalam domain visual guna memfasilitasi identifikasi pola-pola harga yang sering kali berulang dalam periode waktu tertentu. Berbagai jenis representasi grafis, mulai dari grafik garis hingga \textit{candlestick}, digunakan untuk mereduksi derau (\textit{price noise}) harga yang acak menjadi bentuk visual yang lebih intuitif bagi para analis pasar. Prinsip utama dari metodologi ini didasarkan pada postulat bahwa harga pasar saat ini telah mencerminkan secara sempurna seluruh informasi historis serta aksi kolektif para pelaku pasar di masa lampau. Visualisasi data ini sangat membantu agen ekonomi dalam melakukan ekstrapolasi tren masa depan berdasarkan pengenalan pola-pola geometris yang telah terbukti memiliki signifikansi statistik dalam sejarah perdagangan. Kemampuan untuk membaca representasi visual ini menjadi keterampilan dasar bagi praktisi pasar dalam menentukan titik masuk dan keluar yang optimal dalam aktivitas perdagangan.

Pola harga yang terbentuk secara visual dalam grafik teknikal sebenarnya merupakan manifestasi dari psikologi massa serta perilaku kolektif yang mendasari dinamika interaksi antar agen di pasar finansial. Konsep teknikal seperti level \textit{support} dan \textit{resistance} mencerminkan titik-titik memori sosial di mana agen cenderung melakukan reaksi psikologis yang serupa terhadap ambang batas harga tertentu. Transisi dari data harga diskrit ke dalam kerangka waktu kontinu memungkinkan pemodelan kecenderungan sosial tersebut secara matematis melalui pendekatan sosiopsikologi yang terintegrasi dengan hukum fisika statistik. Dengan demikian, analisis teknikal berfungsi sebagai instrumen strategis untuk mendeteksi fenomena pergerakan kawanan (\textit{herding behavior}) yang sering kali menjadi pemicu utama terjadinya volatilitas harga yang ekstrem. Analisis ini memberikan wawasan mengenai sentimen dominan pasar yang dapat digunakan sebagai parameter kontrol dalam model simulasi dinamika harga aset.

\subsection{Analisis Kuantitatif Sebagai Dasar Ekonofisika}
Keuangan kuantitatif (\textit{quantitative finance}) dalam bidang ekonofisika dimulai dengan proses transformasi harga mentah menjadi variabel imbal hasil (\textit{returns}) guna meminimalkan efek non-stasioneritas yang sering kali menghambat analisis statistik. Transformasi ini bertujuan untuk menghasilkan deret waktu yang bersifat stasioner, sehingga berbagai instrumen matematika dapat diimplementasikan secara valid untuk memodelkan perilaku harga di masa depan. Penggunaan imbal hasil logaritmik (\textit{log return}) lebih diutamakan oleh para peneliti karena sifatnya yang memungkinkan agregasi temporal secara linear serta memberikan kemudahan dalam komputasi numerik yang kompleks. Struktur informasi dari data stokastik ini kemudian dianalisis secara mendalam untuk mendekati distribusi normal, yang merupakan prasyarat krusial bagi efisiensi algoritma dalam kerangka teori informasi ekonomi. Pendekatan kuantitatif ini menjembatani antara data pasar mentah yang kacau dengan model fisik yang memiliki tingkat keteraturan matematis yang tinggi.

Proses pengolahan data kuantitatif melibatkan pengambilan statistik imbal hasil serta teknik normalisasi rentang nilai input untuk memastikan konsistensi data sebelum dimasukkan ke dalam model simulasi. Fungsi aktivasi non-linear seperti \textit{Sigmoid} atau \textit{Hyperbolic Tangent} ($tanh$) sering kali diterapkan guna memodelkan fenomena saturasi sentimen pasar serta perilaku agen yang tidak proporsional terhadap perubahan harga. Normalisasi data tersebut sangat penting untuk menjamin bahwa input finansial yang memiliki rentang luas dapat dipetakan secara akurat ke dalam parameter sistem fisik atau sirkuit kuantum. Integrasi antara analisis statistik dan fungsi aktivasi ini memungkinkan model untuk mendeteksi transisi fase pasar, seperti kondisi \textit{crash} atau \textit{rally}, dengan tingkat akurasi yang lebih tinggi dan terukur secara empiris. Penerapan teknik normalisasi ini juga mengurangi beban komputasi pada perangkat keras kuantum dengan memastikan stabilitas nilai parameter selama proses iterasi.

\subsection{Formulasi Objektif Markowitz pada Portofolio Biner}
Strategi diversifikasi fundamental berdasarkan model Markowitz bertujuan untuk mencari titik keseimbangan optimal antara ekspektasi imbal hasil dan tingkat risiko melalui seleksi aset yang paling efisien secara statistik. Dalam konteks portofolio biner, investor dihadapkan pada keputusan mutlak untuk menentukan apakah suatu aset tertentu akan dimasukkan atau dikeluarkan dari keranjang investasi secara definitif. Pergeseran ke domain biner ini mengubah masalah optimasi portofolio konvensional menjadi kategori masalah kombinasional yang bersifat \textit{NP-hard}, sehingga membutuhkan pendekatan komputasi yang lebih canggih dan terspesialisasi. Tantangan utama dalam model ini terletak pada ruang solusi yang tumbuh secara eksponensial seiring dengan bertambahnya jumlah aset, yang menuntut efisiensi tinggi dalam setiap langkah algoritma pencarian solusi. Optimasi portofolio biner ini sangat relevan bagi manajemen aset modern yang membutuhkan keputusan alokasi kapital yang tegas dan bebas dari ambiguitas operasional.

Fungsi biaya dalam model Markowitz menggabungkan komponen matriks kovarians dan ekspektasi imbal hasil melalui parameter penghindaran risiko (\textit{risk aversion}) yang disimbolkan dengan variabel $\lambda$. Matriks kovarians merepresentasikan interaksi risiko sistemik antar aset yang menentukan derajat volatilitas kolektif dalam suatu portofolio, yang merupakan inti dari perhitungan risiko dalam manajemen aset modern. Struktur fungsi ini memiliki karakteristik kuadratik di mana setiap interaksi antar pasangan aset dihitung berdasarkan korelasi pergerakan harga historis yang telah dinormalisasi secara tepat. Formulasi Lagrangian untuk fungsi biaya pada portofolio biner dapat dinyatakan dalam bentuk persamaan berikut:
\begin{equation}
\mathcal{L}(\mathbf{x}) = \frac{\lambda}{2} \sum_{i,j} \sigma_{ij} x_i x_j - \sum_i \mu_i x_i
\end{equation}
di mana $x_i \in \{0,1\}$ merepresentasikan keputusan biner untuk setiap aset $i$, $\sigma_{ij}$ adalah elemen matriks kovarians, dan $\mu_i$ merupakan ekspektasi imbal hasil aset. Pencarian portofolio yang optimal secara teknis setara dengan upaya untuk menemukan kombinasi aset yang mampu meminimalkan risiko total pada target tingkat keuntungan tertentu.

\subsection{Model Mean-Variance Portfolio Optimization}
Model \textit{Mean-Variance Portfolio Optimization} bekerja dengan mengidentifikasi kumpulan solusi terbaik yang dikenal sebagai titik-titik \textit{Pareto Optimum} dalam ruang dimensi antara risiko dan ekspektasi imbal hasil. Titik-titik optimal tersebut kemudian membentuk kurva \textit{Efficient Frontier}, yang berfungsi sebagai acuan standar bagi investor rasional untuk memaksimalkan utilitas investasi mereka di tengah ketidakpastian pasar. Melalui pemanfaatan kurva efisiensi ini, investor dapat secara presisi memilih kombinasi aset yang menawarkan risiko terendah untuk setiap target keuntungan yang telah ditetapkan sebelumnya oleh kebijakan investasi. Implementasi model ini pada aset biner memaksa setiap titik koordinat pada kurva tersebut untuk merepresentasikan keputusan \textit{buy} atau \textit{sell} yang bersifat tegas bagi agen. Kurva efisiensi memberikan batas teoretis yang memisahkan antara strategi investasi yang inferior dengan strategi yang optimal secara statistik.

Bagi para investor institusional, kurva \textit{Efficient Frontier} bertindak sebagai batas atas performa (\textit{benchmark}) yang dapat dicapai secara teoretis dalam kondisi pasar yang berfluktuasi secara dinamis. Pemilihan titik spesifik pada kurva efisiensi ini sangat bergantung pada mandat investasi serta profil risiko yang dimiliki oleh lembaga pengelola dana tersebut dalam jangka panjang. Penggunaan model biner secara inheren memperumit navigasi di sepanjang kurva efisiensi karena setiap pergeseran posisi memerlukan restrukturisasi aset secara total yang berdampak pada biaya transaksi dan likuiditas. Oleh karena itu, akurasi dalam menentukan titik optimal pada model optimasi biner menjadi sangat krusial guna menjaga stabilitas portofolio serta keberlanjutan pertumbuhan kekayaan riil di masa depan. Dinamika ini menuntut pengembangan algoritma optimasi yang mampu memberikan solusi presisi dalam waktu yang efisien guna merespons perubahan kondisi pasar yang cepat.
